\BeleskeID{08-s014}
% element: 08-s014:e03
% element: 08-s014:notes
Izbor manje margine omogućava da se jedan podatak o otpornosti na smetnje koristi i kada nije poznato izlazno logičko stanje. Ako bismo izabrali veću marginu, smetnja tog iznosa ne bi nužno bila dozvoljena i za drugo stanje.

Sabiranjem dve razlike napona i grupisanjem članova dobija se
\[NM_L+NM_H=(V_{IL}-V_{OL})+(V_{OH}-V_{IH})=(V_{OH}-V_{OL})-(V_{IH}-V_{IL}).\]
Zbir margina je, dakle, izlazni naponski raspon umanjen za širinu prelaznog ulaznog opsega. Za isti izlazni raspon širi prelazni opseg ostavlja manji zbir margina. Pri prikazanom poretku svih nivoa unutar napojnog opsega izlazni raspon nije veći od $V_{CC}$, pa ni zbir margina ne može biti veći od njega. To objašnjava proveru iz ispitnih zadataka.
